\documentclass[11pt,letterpaper]{article}
\usepackage[T1]{fontenc}
\usepackage[utf8]{inputenc}
\usepackage{lmodern,microtype}
\usepackage[margin=0.78in,headheight=15pt]{geometry}
\usepackage{amsmath,amssymb,booktabs,longtable,tabularx,array,enumitem,xcolor,hyperref,fancyhdr,lastpage,needspace}
\definecolor{EngineBlue}{HTML}{183B56}
\definecolor{EngineTeal}{HTML}{117A83}
\definecolor{EngineGray}{HTML}{596671}
\definecolor{EngineLight}{HTML}{EEF4F7}
\hypersetup{colorlinks=true,urlcolor=EngineTeal,linkcolor=EngineBlue}
\setlength{\parindent}{0pt}
\setlength{\parskip}{5pt}
\emergencystretch=3em
\setlist{nosep,leftmargin=*}
\pagestyle{fancy}\fancyhf{}
\fancyhead[L]{\small\sffamily Game Engine Architecture | Engineering Suite}
\fancyhead[R]{\small\sffamily \DocType}
\fancyfoot[L]{\scriptsize Source-grounded engineering curriculum | v1.0}
\fancyfoot[R]{\scriptsize \thepage/\pageref{LastPage}}
\newcommand{\DocType}{Reference}
\newcommand{\SuiteHeading}[3]{%
\renewcommand{\DocType}{#2}%
{\small\bfseries\sffamily\color{EngineTeal}#2\hfill #3}\par\vspace{6pt}
{\LARGE\bfseries\color{EngineBlue}#1}\par\vspace{4pt}
{\color{EngineTeal}\rule{\linewidth}{1pt}}\par\vspace{7pt}}
\newcommand{\LabEvidence}[1]{\par\textbf{Evidence required:} #1\par}
\newcommand{\WriteLines}[1]{\par\begingroup\setlength{\parskip}{9pt}\count255=0\loop\ifnum\count255<#1\noindent\rule{\linewidth}{0.35pt}\par\advance\count255 by1\repeat\endgroup\smallskip}
\newcommand{\WritingBox}[1]{\noindent\fbox{\parbox[t][#1][t]{0.97\linewidth}{\mbox{}}}\par\smallskip}
\newcommand{\SourceNote}{\footnotesize\textit{Primary reading: Jason Gregory, \emph{Game Engine Architecture}, 3rd ed. (CRC Press, 2018). The narrative, exercises, laboratory designs and scoring instruments in this suite are independently authored teaching materials, not a reproduction of the textbook.}\normalsize\par}
\newcommand{\Criterion}[2]{\textbf{#1.} #2\par}
\newcommand{\FormField}[2]{\Needspace{3\baselineskip}\textbf{#1}\par\textcolor{EngineGray}{\footnotesize #2}\par\WriteLines{2}}
\newcommand{\BlankField}[1]{\Needspace{3\baselineskip}\textbf{#1}\par\WriteLines{2}}
\newcolumntype{Y}{>{\raggedright\arraybackslash}X}

\begin{document}
\SuiteHeading{Audio and Input Integration}{Technical reference}{Runtime systems}
\SourceNote
\section{Scope and normative conventions}
This reference summarizes design options and independently proposed engineering controls. It is not a statement that the textbook mandates any particular API or modern implementation choice.

\subsection{Input ingestion}
Raw device events are timestamped, normalized and mapped to semantic actions; recorded traces can drive deterministic simulation testing.

\begin{itemize}
\item Define the observable interface and resource owner.
\item Specify a success invariant and a deliberate failure case.
\item Trace at least one update frame from input to output.
\end{itemize}
\subsection{Audio mixing}
Sample-rate, channel layout, gain and buffer size are ABI-like constraints. Sum voices with headroom and validate clipping.

\begin{itemize}
\item Define the observable interface and resource owner.
\item Specify a success invariant and a deliberate failure case.
\item Trace at least one update frame from input to output.
\end{itemize}
\subsection{Realtime boundaries}
Audio callbacks must not block on asset I/O, general-purpose locks, allocation or console output; push prevalidated commands over a bounded queue.

\begin{itemize}
\item Define the observable interface and resource owner.
\item Specify a success invariant and a deliberate failure case.
\item Trace at least one update frame from input to output.
\end{itemize}
\subsection{Spatial sound}
Attenuation and panning are approximation models; listener coordinate transform must match renderer/physics conventions.

\begin{itemize}
\item Define the observable interface and resource owner.
\item Specify a success invariant and a deliberate failure case.
\item Trace at least one update frame from input to output.
\end{itemize}
\section{Design review questions}
\begin{itemize}
\item Where are time, memory and object lifetimes transferred across subsystem boundaries?
\item Which statements are assumptions, and what test would invalidate them?
\item What fails if content is unloaded while a worker thread still holds a reference?
\item Which metrics can establish that this subsystem satisfies a frame-time budget?
\end{itemize}
\section{Evidence package}
\begin{itemize}
\item Versioned contract and small dependency diagram.
\item Tests covering invalid inputs and state transitions.
\item Reproducible benchmark data, limitations and release decision.
\end{itemize}

\end{document}
